%% ====================================================================
%%  Salt Pans of India: A Comprehensive Study
%%  LaTeX Source (compile with: pdflatex -> bibtex -> pdflatex x2)
%%  Contributors: Arnav Gupta · Satvik Srinivas
%%  Date: September 2026
%% ====================================================================
\documentclass[12pt,a4paper]{article}

%% ── Packages ──────────────────────────────────────────────────────────
\usepackage[a4paper, margin=2.5cm]{geometry}
\usepackage[utf8]{inputenc}
\usepackage[T1]{fontenc}
\usepackage{lmodern}
\usepackage{microtype}
\usepackage{setspace}
\usepackage{parskip}
\usepackage{amsmath,amssymb}
\usepackage{booktabs}
\usepackage{longtable}
\usepackage{tabularx}
\usepackage{array}
\usepackage{multirow}
\usepackage{xcolor}
\usepackage{fancyhdr}
\usepackage{titlesec}
\usepackage{tcolorbox}
\usepackage{enumitem}
\usepackage{hyperref}
\usepackage{natbib}
\usepackage{graphicx}
\usepackage{float}
\usepackage{caption}

%% ── Colors ────────────────────────────────────────────────────────────
\definecolor{deepteal}{HTML}{00524A}
\definecolor{medteal}{HTML}{007A6E}
\definecolor{lightteal}{HTML}{E0F4F2}
\definecolor{accent}{HTML}{C8820A}
\definecolor{muted}{HTML}{5A5A5A}
\definecolor{bodytext}{HTML}{2D2D2D}

%% ── Hyperref ──────────────────────────────────────────────────────────
\hypersetup{
  colorlinks=true,
  linkcolor=deepteal,
  citecolor=medteal,
  urlcolor=accent,
  pdftitle={Salt Pans of India: A Comprehensive Study},
  pdfauthor={Arnav Gupta; Satvik Srinivas},
  pdfsubject={Indian Salt Industry, Worker Welfare, Quality Improvement}
}

%% ── Typography ────────────────────────────────────────────────────────
\setstretch{1.2}
\setlength{\parindent}{0pt}
\setlength{\parskip}{6pt}

%% ── Section Formatting ────────────────────────────────────────────────
\titleformat{\section}
  {\color{white}\Large\bfseries\sffamily}
  {\colorbox{deepteal}{\parbox[c][1.4em]{0.9em}{\centering\thesection}}}
  {0.6em}
  {\colorbox{deepteal}{\parbox[t]{\dimexpr\linewidth-2.4em}{\vspace{2pt}#1\vspace{2pt}}}}

\titleformat{\subsection}
  {\color{medteal}\large\bfseries\sffamily}
  {\thesubsection}
  {0.5em}
  {}

\titleformat{\subsubsection}
  {\color{accent}\normalsize\bfseries\itshape\sffamily}
  {\thesubsubsection}
  {0.4em}
  {}

\titlespacing{\section}{0pt}{18pt}{8pt}
\titlespacing{\subsection}{0pt}{12pt}{4pt}
\titlespacing{\subsubsection}{0pt}{10pt}{3pt}

%% ── Page Style ────────────────────────────────────────────────────────
\pagestyle{fancy}
\fancyhf{}
\fancyhead[L]{\color{muted}\small\sffamily Salt Pans of India: A Comprehensive Study}
\fancyhead[R]{\color{muted}\small\sffamily Gupta \& Srinivas, 2026}
\fancyfoot[C]{\color{muted}\small---~\thepage~---}
\renewcommand{\headrulewidth}{0.5pt}
\renewcommand{\headrule}{\color{accent}\hrule height\headrulewidth width\headwidth}
\renewcommand{\footrulewidth}{0pt}

%% ── tcolorbox Styles ─────────────────────────────────────────────────
\tcbuselibrary{skins,breakable}

\newtcolorbox{abstractbox}{
  enhanced, breakable,
  colback=lightteal, colframe=deepteal,
  boxrule=1.2pt, arc=6pt,
  left=8pt, right=8pt, top=8pt, bottom=8pt,
  borderline west={4pt}{0pt}{accent}
}

\newtcolorbox{infobox}[1][]{
  enhanced,
  colback=lightteal!60, colframe=medteal,
  boxrule=0.8pt, arc=4pt,
  fonttitle=\bfseries\sffamily\color{white},
  attach boxed title to top left={yshift=-2mm,xshift=4mm},
  boxed title style={colback=medteal, boxrule=0pt, arc=3pt},
  #1
}

%% ── Table column types ────────────────────────────────────────────────
\newcolumntype{L}[1]{>{\raggedright\arraybackslash}p{#1}}
\newcolumntype{C}[1]{>{\centering\arraybackslash}p{#1}}
\newcolumntype{R}[1]{>{\raggedleft\arraybackslash}p{#1}}

%% ── Caption style ─────────────────────────────────────────────────────
\captionsetup{font={small,it}, labelfont={bf,color=medteal}}

%% ====================================================================
%%  DOCUMENT
%% ====================================================================
\begin{document}

%% ── Title Page ────────────────────────────────────────────────────────
\begin{titlepage}
\begin{center}
  \vspace*{2cm}
  {\fontsize{32}{40}\selectfont\bfseries\sffamily\color{deepteal}
    Salt Pans of India}\\[0.8em]
  {\Large\sffamily\color{medteal}
    A Comprehensive Study of Production, Worker Welfare,\\[0.3em]
    Quality Improvement, and Ecological Significance}

  \vspace{1.5em}
  {\color{accent}\rule{0.6\textwidth}{2pt}}

  \vspace{1.5em}
  {\large\bfseries\sffamily Arnav Gupta \quad $\cdot$ \quad Satvik Srinivas}\\[0.4em]
  {\itshape\color{muted} Lavana Project Research Initiative}\\[0.3em]
  {\color{muted} September 2026}

  \vspace{2.5em}
  \begin{abstractbox}
    {\bfseries\sffamily\color{deepteal} Abstract}\\[0.5em]
    India is the world's third-largest producer of salt, generating approximately
    \textbf{266 lakh tonnes} (26.6 million tonnes) annually, of which \textbf{Gujarat
    alone contributes 85.6\%}. Salt pans—ranging from vast coastal solar-evaporation
    works and inland saline lakes to traditional community pans—underpin a sector that
    employs over 77,000 registered workers and sustains an informal workforce estimated
    at several hundred thousand. This report synthesises geographic, socioeconomic, and
    ecological data drawn from the Lavana GIS dataset (48 sites across 11 states/UTs),
    the Indian Minerals Yearbook 2022, the Salt Commissioner's Organisation, peer-reviewed
    literature, and field journalism to present a holistic account of India's salt pans.
    Topics covered include: the geography and geology of major production centres;
    historical significance; working conditions, health challenges, and welfare initiatives
    targeting the salt-pan labour force; best practices in quality improvement and
    iodisation; technological modernisation; ecological roles in supporting migratory
    bird populations; and the regulatory and legal landscape governing this ancient industry.
  \end{abstractbox}

  \vspace{1.2em}
  {\small\itshape\color{muted}
    \textbf{Keywords:} salt pans, India, Agariyas, solar evaporation, iodisation,
    worker welfare, flamingos, Little Rann of Kutch, Sambhar Lake, Lavana GIS dataset}
\end{center}
\end{titlepage}

\tableofcontents
\clearpage

%% ====================================================================
\section{Introduction}
%% ====================================================================

Salt — chemically sodium chloride (NaCl) — is among humanity's oldest and most
essential commodities. In India, the salt industry carries profound historical weight,
having been at the centre of Mahatma Gandhi's celebrated \textbf{Salt Satyagraha
(Dandi March, 1930)}, a defining act of civil disobedience against British colonial
taxation. Today, India ranks third globally in salt production (after China and the
United States), producing roughly \textbf{30--32 million tonnes per year} in the
2022--2024 period.

The industry is structured around three broad production systems:
\begin{enumerate}[leftmargin=*, label=(\roman*)]
  \item \textbf{Coastal solar-evaporation works}, which dominate Gujarat, Tamil Nadu,
        Andhra Pradesh, Maharashtra, and Goa;
  \item \textbf{Inland saline lake and sub-soil brine operations}, centred on
        Rajasthan's lakes (Sambhar, Pachpadra, Didwana) and Gujarat's Little Rann
        of Kutch; and
  \item \textbf{Smaller artisanal and cooperative pans} found along the Konkan,
        Odisha, and West Bengal coasts.
\end{enumerate}

Despite its scale and strategic importance, the industry is characterised by informality,
low mechanisation, and persistent deficits in worker welfare. This report aims to provide
a data-grounded, multi-dimensional analysis of the sector to inform policy, research,
and public understanding.

\subsection{Scope and Data Sources}

The primary geographic dataset used in this report is the \textbf{Lavana GIS Dataset v1}
(compiled September 2026), which documents 48 salt pan sites across 11 Indian states and
union territories. Quantitative production statistics are drawn from the
\textbf{Indian Minerals Yearbook 2022} (IBM/Salt Commissioner's Organisation).
Socioeconomic data are sourced from the Journal of Management Research and Analysis
(JMRA, 2025), the Berkeley Journal of Sociology, PARI (People's Archive of Rural India),
Mongabay India, and The Hindu. Ecological information is drawn from the Ramsar Site
Information Service (RSIS), Bombay Natural History Society (BNHS), and peer-reviewed
ornithological literature.

%% ====================================================================
\section{Geography and Production Landscape}
%% ====================================================================

India's coastline stretches over 7,500 km along the Arabian Sea, Bay of Bengal, and
Indian Ocean, providing extensive evaporation potential. The interior plateaux of
Rajasthan add significant inland saline resources. Salt production is highly
concentrated: \textbf{Gujarat accounts for over 85\%} of national output.

\subsection{State-wise Production (2021--22)}

\begin{table}[H]
\centering
\caption{State-wise salt production, 2021--22.
Source: IBM Indian Minerals Yearbook 2022 / Salt Commissioner's Organisation.
* AP figure from \textit{Hans India} (Apr 2022); \dag{} journalist estimate (PARI);
\ddag{} 2008 district figure; \S{} manufacturer self-report (\textit{Mongabay}, 2024).}
\smallskip
\begin{tabular}{L{4.5cm} C{2.8cm} C{1.8cm} L{5cm}}
\toprule
\rowcolor{deepteal}
\color{white}\textbf{State / UT} &
\color{white}\textbf{Production (Lakh t)} &
\color{white}\textbf{Share (\%)} &
\color{white}\textbf{Primary Method} \\
\midrule
Gujarat & 227.64 & 85.6 & Coastal solar / sub-soil brine \\
\rowcolor{lightteal}
Tamil Nadu & 17.21 & 6.5 & Coastal solar evaporation \\
Rajasthan & 16.90 & 6.4 & Lake brine / sub-soil brine \\
\rowcolor{lightteal}
Andhra Pradesh & 4.35* & 1.6 & Coastal solar evaporation \\
Maharashtra & --- & $\sim$0.5 & Coastal solar evaporation \\
\rowcolor{lightteal}
Other States & $\sim$4.24 & 1.6 & Various \\
\midrule
\textbf{India Total} & \textbf{266.00} & \textbf{100} & --- \\
\bottomrule
\end{tabular}
\end{table}

\subsection{Gujarat: The Dominant Producer}

Gujarat's dominance stems from its 1,600 km coastline, arid climate, low annual
rainfall (400--600 mm on the Saurashtra coast), and high solar insolation. The
\textbf{Little Rann of Kutch (LRK)}—a seasonal saline marsh of $\sim$5,000 km$^2$—is
home to the \textit{Agariya} community, traditional sub-soil brine salt farmers.
Approximately \textbf{45,000 Agariyas} work within or adjacent to the Indian Wild Ass
Sanctuary, producing the iconic ``white gold'' of Kutch.

Average yields have declined: from approximately \textbf{80 t/acre in 2016--17 to
68 t/acre in 2020--21} in Jamnagar (JMRA, 2025), driven by erratic monsoon patterns
and increasingly frequent cyclones that disrupt the October--May production season.

\subsection{Rajasthan: Inland Saline Heritage}

\textbf{Sambhar Salt Lake} (Ramsar site No.\ 464, designated 1990) is India's largest
inland saline lake with a catchment of $\sim$5,700 km$^2$. Sambhar Salts Ltd produces
approximately \textbf{2 lakh tonnes/yr} of raw salt. \textbf{Pachpadra Lake} near
Balotra has been mined by the traditional Khara (Kharwal) community for centuries using
the \textit{morli} shrub crystallisation technique; its salt is approximately 98\% NaCl.
\textbf{Didwana Lake} produces lower-grade industrial salt, primarily used for sodium
sulphate extraction.

\subsection{Tamil Nadu, Andhra Pradesh, and Other States}

Tamil Nadu's production belt stretches from Thoothukudi (``Salt Capital'') northward
through Marakkanam and Vedaranyam, producing approximately \textbf{17.21 lakh tonnes/yr}.
Andhra Pradesh's coast from Nellore to Srikakulam—especially the ``Salt Bowl'' of
Naupada—contributes roughly 4.35 lakh tonnes/yr. Maharashtra's Mumbai salt pans
($\sim$5,379 acres) serve a dual ecological--industrial function; much of this land is
under development pressure for urban infrastructure.

%% ====================================================================
\section{Historical Significance}
%% ====================================================================

\subsection{Ancient and Medieval Trade}

Large-scale salt production around Patadi, Jhinjuwada, and Kharaghoda in the Little
Rann of Kutch is documented as early as the \textbf{10th century CE} (Campbell, 1887,
via Berkeley Journal). A \textbf{Mughal firman of 1669--70} reinstating the King of
Halvad as owner of \textit{agar} (salt pans) illustrates the commodity's political
economy in the pre-colonial period. Between \textbf{1870 and 1905--06}, approximately
\textbf{3.67 million tonnes} of salt were extracted from Sambhar alone.

\subsection{The British Salt Monopoly}

Under colonial rule, the 1882 \textbf{Salt Act} established a British monopoly,
criminalising the independent production or collection of salt. This generated enormous
revenue but imposed a regressive burden on ordinary Indians. The Tata group's acquisition
of the Okha Salt Works at Mithapur in \textbf{1939} (with a soda ash plant starting
production in 1944) represents early industrialisation of the sector under private enterprise.

\subsection{The Dandi March (Salt Satyagraha, 1930)}

On \textbf{12 March 1930}, Mahatma Gandhi and 78--80 volunteers began a 241-mile
(387 km) march from Sabarmati Ashram, Ahmedabad, to the coastal village of Dandi,
Gujarat. Upon reaching Dandi on \textbf{6 April 1930}, Gandhi broke the Salt Act by
collecting natural salt from the seashore, triggering nationwide civil disobedience.
Simultaneously, C.\ Rajagopalachari led a parallel salt march from Tiruchirapalli to
\textbf{Vedaranyam}, Tamil Nadu—now a Ramsar-listed site—underscoring that salt
transcended geography in its unifying power.

\subsection{Post-Independence Development}

After 1947, the Salt Commissioner's Organisation (SCO) became the nodal body for
production regulation, land leasing, welfare, and iodisation policy. Universal Salt
Iodisation (USI) became national policy in the 1990s to address iodine-deficiency
disorders; FSSAI now mandates specific iodine concentrations in all edible salt.

%% ====================================================================
\section{The Salt Pan Workforce: Conditions and Welfare}
%% ====================================================================

The salt industry employs approximately \textbf{77,086 average workers} (IBM IMYB 2022,
registered/formal). The informal workforce—including migrant Agariya families in Gujarat
and seasonal workers in Tamil Nadu and Andhra Pradesh—is estimated at several hundred
thousand.

\subsection{Working Conditions}

Salt workers endure some of the harshest occupational conditions in India:

\begin{itemize}[leftmargin=*, nosep]
  \item \textbf{Extreme heat:} Surface temperatures routinely reach 45--50°C;
        reflective white salt amplifies UV radiation.
  \item \textbf{Saline immersion:} Workers wade through hypersaline brine for hours,
        causing rapid skin degradation on feet and lower limbs.
  \item \textbf{Seasonal precarity:} Coastal operations run October--May; during the
        monsoon, income ceases entirely.
  \item \textbf{Lack of PPE:} Most workers lack rubber boots, gloves, and UV-protective
        eyewear.
  \item \textbf{Absence of amenities:} Potable water, sanitation, rest sheds, and
        childcare are absent at most pan sites.
  \item \textbf{Child labour and migration:} Agariya families relocate to the LRK for
        months, often with children who miss schooling.
\end{itemize}

\subsection{Health Impact Assessment}

\begin{table}[H]
\centering
\caption{Occupational health concerns among salt pan workers.
Sources: NIH studies; IJFMR; Human Rights Research (2024); JMRA (2025).}
\smallskip
\begin{tabular}{L{3.5cm} L{6cm} C{2.5cm}}
\toprule
\rowcolor{medteal}
\color{white}\textbf{Health Concern} &
\color{white}\textbf{Cause / Mechanism} &
\color{white}\textbf{Prevalence} \\
\midrule
Skin lesions \& burns & Prolonged saline contact; sun at 45--50°C & Very High \\
\rowcolor{lightteal}
Eye irritation / vision loss & Salt spray; UV; no PPE & High \\
Musculoskeletal disorders & Repetitive harvesting; load carrying & High \\
\rowcolor{lightteal}
Kidney dysfunction & Chronic dehydration; high NaCl intake & Moderate--High \\
Hypertension & High sodium intake; heat stress & Moderate \\
\rowcolor{lightteal}
Respiratory issues & Salt dust during harvesting & Moderate \\
Malnutrition / anaemia & Low wages; seasonal unemployment & High (women) \\
\bottomrule
\end{tabular}
\end{table}

\subsection{Government and NGO Welfare Initiatives}

\subsubsection{Tamil Nadu Salt Pan Workers Welfare Board (2023)}

In \textbf{July 2023}, the Government of Tamil Nadu established a Welfare Board for
Salt Pan Workers under the Tamil Nadu Manual Workers Act, 1982. The board provides:
financial relief during the monsoon; health insurance; housing assistance; and educational
scholarships. As of mid-2026, benefit delivery remains slow (\textit{The Hindu}; Tamil Vahini, 2026).

\subsubsection{Gujarat Salt Empowered Committee}

The Gujarat government's \textbf{Salt Empowered Committee} deploys health vans offering
lab services and maternity testing at LRK pan sites. In partnership with \textbf{SEWA},
the scheme has promoted solar-powered brine pumps, reducing physical labour and diesel
expenditure. HSL (Hindustan Salts Ltd) operates a hospital in Kharaghoda.

\subsubsection{Proposed National Salt Workers Welfare Bill}

A \textbf{Salt Workers Welfare Bill} has been introduced in Parliament in 2014, 2023,
and 2024. The proposed legislation would establish a \textbf{National Commission for the
Welfare of Salt Workers} with powers to register all workers nationally; set minimum
wages; mandate PPE; and ensure regular health check-ups (GeoJuris Today, 2024).

\subsubsection{Agariya Land Rights}

Agariya families work within the \textbf{Indian Wild Ass Sanctuary} (4,953.71 km$^2$),
creating a tension between conservation law and livelihood. A formal land-rights decision
remains pending as of 2025 (Berkeley Journal, 2025).

\subsection{Remaining Gaps and Priority Actions}

\begin{itemize}[leftmargin=*]
  \item Enact and fund the national Salt Workers Welfare Bill.
  \item Extend ESI and Provident Fund to cover informal and migrant salt workers.
  \item Mandate provision of potable water, rest sheds, and sanitation at all pan sites.
  \item Subsidise PPE distribution (boots, gloves, UV goggles) through SCO.
  \item Scale up mobile health units to Andhra Pradesh, Tamil Nadu, and Maharashtra.
  \item Formalise Agariya land tenure under the Forest Rights Act.
\end{itemize}

%% ====================================================================
\section{Quality Improvement and Technological Modernisation}
%% ====================================================================

\subsection{Technology and Best Practices}

\begin{table}[H]
\centering
\caption{Key technologies for quality improvement in Indian salt production.
Sources: CSIR-CSMCRI; SEWA; SCO; FSSAI; \textit{Mongabay India} (2024).}
\smallskip
\begin{tabular}{L{4cm} L{5cm} L{4.5cm}}
\toprule
\rowcolor{medteal}
\color{white}\textbf{Technology / Practice} &
\color{white}\textbf{Benefit} &
\color{white}\textbf{Status in India} \\
\midrule
Geo-Membrane (GM) Liners & Reduces seepage 15--20\%; improves purity 12--15\% &
  Pilot scale; CSIR-CSMCRI \\
\rowcolor{lightteal}
Concrete / fly-ash beds & Prevents contamination; enables mechanical harvesting &
  Partial adoption; larger works \\
Solar PV brine pumps & Replaces diesel; lowers cost \& carbon footprint &
  SEWA-supported rollout, Gujarat \\
\rowcolor{lightteal}
Heat exchanger / solar evaporation & Yield 60$\to$85 t/acre; cycle $-$30--40\% &
  Research stage; CSIR-CSMCRI \\
Mechanised harvesting & Improves efficiency; reduces impurities &
  Large private works; not SMEs \\
\rowcolor{lightteal}
Automated spray iodisation & Uniform iodine mixing; FSSAI compliance &
  Mandated; adoption ongoing \\
Model Salt Farms & Demonstrates best practices to smaller producers &
  SCO initiative; GJ, TN, Odisha \\
\rowcolor{lightteal}
Salt Testing Kits (STKs) & Field-level iodine quality verification &
  Distributed by state health depts. \\
\bottomrule
\end{tabular}
\end{table}

\subsection{Universal Salt Iodisation (USI)}

India's USI programme mandates that all edible salt be iodised at \textbf{15 ppm at
the consumer level} (FSSAI IS 7224:2018). Iodisation is performed using the
\textbf{spray method}, in which a potassium iodate solution is uniformly applied
during packaging. Quality control involves: SCO sampling at production sites;
distribution of Salt Testing Kits (STKs) to state health directorates; mandatory
iodine content labelling; and prosecution of producers failing FSSAI standards.

\subsection{Climate Impacts on Quality and Yield}

Climate change poses a dual threat: erratic monsoon onset/withdrawal shortens the
production window, while cyclone frequency has increased. Manufacturers in Bhavnagar
report \textbf{25\% production losses} in years with two cyclone events
(\textit{Mongabay}, 2024). Adaptation strategies include:
\begin{itemize}[leftmargin=*]
  \item Geo-membrane liners in concentrator ponds to prevent dilution and seepage.
  \item Weather-based insurance pilots for salt workers (proposed by SCO).
  \item Early-warning integration with the India Meteorological Department.
\end{itemize}

\subsection{Proposed National Institute for Salt Technology}

The Marine Salt Manufacturers Association has advocated for a dedicated
\textbf{National Institute for Salt Technology} to conduct applied R\&D on evaporation
science and crystallisation chemistry; provide skill development and certification;
and benchmark Indian salt quality against global standards for export competitiveness.
No such institution exists as of 2026.

%% ====================================================================
\section{Ecological Significance}
%% ====================================================================

Salt pans — though largely human-made — function as critical wetland ecosystems under
the Ramsar Convention framework and India's Wetlands (Conservation and Management)
Rules 2017. Their shallow, hypersaline waters and exposed mudflats provide irreplaceable
habitat for migratory and resident waterbirds along the \textbf{Central Asian Flyway}.

\subsection{Bird Habitat and Migratory Importance}

Key documented observations include:
\begin{itemize}[leftmargin=*]
  \item \textbf{Thane Creek Flamingo Sanctuary} (Mumbai, Ramsar Aug.\ 2022): BNHS
        recorded over \textbf{130,000 flamingos} in 2022; adjacent salt pan areas
        form essential foraging grounds.
  \item \textbf{Great Rann of Kutch / LRK}: breeding ground for Greater Flamingo;
        seasonal water bodies support Indian Wild Ass and thousands of waders.
  \item \textbf{Sambhar Lake} (Ramsar \#464): wintering and staging ground; a 2019
        avian botulism event killed thousands of birds.
  \item \textbf{Vedaranyam / Point Calimere} (Ramsar \#1210): $\sim$257 bird species;
        up to \textbf{30,000 flamingos} in peak season; Spoonbill Sandpiper documented.
  \item \textbf{Tata Chemicals Mithapur}: biodiversity programme documents over
        \textbf{150 bird species} within the industrial works estate.
\end{itemize}

\subsection{Ramsar and Protected Area Status}

Of the 43 sites in the Lavana Dataset v1, \textbf{2 are confirmed Ramsar sites} and
\textbf{3 are within designated Protected Areas}. Many more overlap with Important
Bird Areas (IBAs) and Key Biodiversity Areas (KBAs), though formal designation has
not been extended to most.

\subsection{Urban Salt Pans as Flood Buffers}

The salt pans of Mumbai's eastern suburbs provide an often-overlooked hydrological
service: \textbf{absorption of monsoon floodwater}. During the 2005 Mumbai floods,
intact salt pan areas demonstrably reduced inundation in adjacent neighbourhoods.
The redevelopment of 256 acres of Mumbai salt pans for the Dharavi rehabilitation
project (approved 2024) has raised significant concern among urban ecologists.

\subsection{Conservation Challenges}

\begin{itemize}[leftmargin=*]
  \item \textbf{Urban encroachment:} Development pressure on periurban pans reduces
        habitat extent.
  \item \textbf{Unregulated borewell extraction:} At Sambhar Lake, illegal pumping
        reduces lake water levels and alters salinity gradients.
  \item \textbf{Pollution:} Pesticide runoff from agriculture and shrimp farms
        degrades water quality at Point Calimere and nearby pans.
  \item \textbf{Sand mining:} Illegal sand extraction at Chinnaganjam (AP) destroys
        pan infrastructure.
  \item \textbf{Avian disease:} The 2019 Sambhar botulism event demonstrated
        catastrophic mortality risk at saline aggregation sites.
\end{itemize}

%% ====================================================================
\section{Regulatory and Legal Framework}
%% ====================================================================

\begin{itemize}[leftmargin=*]
  \item \textbf{Salt Commissioner's Organisation (SCO):} Central authority (under DPIIT)
        responsible for production regulation, quality control, iodisation oversight,
        land leasing, and welfare schemes.
  \item \textbf{Coastal Regulation Zone (CRZ) Rules:} Salt pans in coastal areas are
        classified as CRZ-I(B), permitting only salt extraction and natural gas exploration.
  \item \textbf{Wildlife Protection Act 1972:} Governs the LRK/Indian Wild Ass
        Sanctuary, creating land-use restrictions that conflict with Agariya livelihoods.
  \item \textbf{Ramsar Convention:} Salt pans with significant bird populations but
        without formal designation lack equivalent international conservation status.
  \item \textbf{FSSAI Food Safety Standards:} Mandates iodine content in edible salt
        and regulates packaging and labelling.
  \item \textbf{Forest Rights Act 2006:} Potentially applicable to Agariya customary
        rights within the sanctuary.
\end{itemize}

%% ====================================================================
\section{Best Practices and Policy Recommendations}
%% ====================================================================

\subsection{Production and Quality}

\begin{itemize}[leftmargin=*]
  \item Mainstream geo-membrane liner technology through SCO-subsidised loans to
        small and medium salt manufacturers.
  \item Establish Model Salt Farms in every salt-producing state.
  \item Create a National Institute for Salt Technology for R\&D and quality benchmarking.
  \item Integrate real-time IMD weather alerts with pan management protocols.
  \item Expand solar PV brine pump subsidies to all states.
\end{itemize}

\subsection{Worker Welfare}

\begin{itemize}[leftmargin=*]
  \item Urgently enact the Salt Workers Welfare Bill to create a national framework.
  \item Extend ESI and EPFO coverage to all salt workers, including seasonal migrants.
  \item Mandate PPE provision by pan operators as a licensing condition.
  \item Require all salt works to provide potable water, sanitation, rest sheds,
        and crèche facilities.
  \item Scale mobile health van services to all major salt-producing districts.
  \item Formalise Agariya land tenure under Forest Rights Act provisions.
\end{itemize}

\subsection{Ecological Conservation}

\begin{itemize}[leftmargin=*]
  \item Designate ecologically significant salt pans as Key Biodiversity Areas.
  \item Establish buffer zones around Ramsar-designated saline lakes to restrict
        unregulated brine extraction.
  \item Implement a Sambhar Lake Restoration Plan.
  \item Adopt a ``Working Wetlands'' framework that formalises the dual
        productive--ecological role of human-made salt pans.
  \item Retain Mumbai salt pan CRZ-I(B) status.
\end{itemize}

%% ====================================================================
\section{Conclusion}
%% ====================================================================

India's salt pan landscape is a palimpsest of history, ecology, and human endeavour.
From the ancient sub-soil brine fields of the Little Rann of Kutch to the gleaming
coastal pans of Thoothukudi, from Sambhar's flamingo-filled shallows to Mumbai's
threatened urban wetlands, salt pans are far more than industrial sites—they are living
heritage landscapes that sustain livelihoods, support migratory biodiversity, and buffer
coastal communities against floods.

The sector faces a convergence of stressors: climate volatility, development pressure,
persistent worker welfare deficits, and the absence of a coherent national legislative
framework for the informal workforce. The good news is that solutions are well-understood.
Geo-membrane technology, solar pumps, automated iodisation, and mobile health services
have been piloted successfully. The Tamil Nadu Welfare Board, SEWA's Gujarat programmes,
and SCO's Model Salt Farm initiative demonstrate what coordinated action can achieve.

What is needed is scale, political will, and institutional architecture: a national Salt
Workers Welfare Act, a dedicated technology institute, and a formal ``Working Wetlands''
policy that recognises the ecological services salt pans provide. India's third-place
ranking in global salt production should be matched by first-place commitment to the human
dignity of those who produce it and the ecological systems it sustains.

%% ── Acknowledgements ─────────────────────────────────────────────────
\medskip
\noindent\textbf{Acknowledgements}

The authors gratefully acknowledge the open data resources of the Salt Commissioner's
Organisation (SCO/DPIIT), the Indian Bureau of Mines (IBM), the Ramsar Sites Information
Service (RSIS), the People's Archive of Rural India (PARI), Mongabay India, The Hindu,
the Berkeley Journal of Sociology, and the Bombay Natural History Society (BNHS).
Special thanks to the Agariya communities of the Little Rann of Kutch and the salt
workers of Tamil Nadu and Andhra Pradesh whose labour and knowledge underpin this industry.

%% ====================================================================
%%  REFERENCES
%% ====================================================================
\clearpage
\section*{References}
\addcontentsline{toc}{section}{References}

\begin{enumerate}[leftmargin=*, label={[\arabic*]}]
  \item Indian Bureau of Mines (2022). \textit{Indian Minerals Yearbook 2022 --- Salt}.
        IBM, Government of India. \url{https://ibm.gov.in/}

  \item Salt Commissioner's Organisation (2024). \textit{Salient Features of Indian Salt
        Industry}. DPIIT, Ministry of Commerce and Industry.
        \url{https://saltcomindia.gov.in/}

  \item Chaudhary, P.\ et al.\ (2025). Performance evaluation of the salt industry in
        India with special reference to Gujarat. \textit{Journal of Management Research
        and Analysis}, 11(4). \url{https://jmra.in/archive/volume/11/issue/4/article/16057}

  \item Berkeley Journal of Sociology (2025). Salt `Farming': Gendered Labour and
        Ecologies of migrant workers in Little Rann of Kutch.
        \url{https://berkeleyjournal.org/2025/06/11/salt-farming/}

  \item People's Archive of Rural India (PARI). The Rani of Thoothukudi's salt pans.
        \url{https://ruralindiaonline.org/article/the-rani-of-thoothukudis-salt-pans}

  \item Mongabay India (2024, March). Uncertain weather makes that pinch of salt dearer.
        \url{https://india.mongabay.com/2024/03/uncertain-weather-makes-that-pinch-of-salt-dearer/}

  \item Ramsar Sites Information Service (RSIS). Sambhar Lake --- Ramsar site no.\ 464
        (designated 23 March 1990). \url{https://rsis.ramsar.org/ris/464}

  \item Ramsar Sites Information Service (RSIS). Point Calimere Wildlife and Bird
        Sanctuary --- Ramsar site no.\ 1210 (designated 19 August 2002).
        \url{https://rsis.ramsar.org/ris/1210}

  \item Environment \& Society Portal (2026). Between Salt and Water: The Environmental
        Crisis at Sambhar Lake, Rajasthan, India.
        \url{https://www.environmentandsociety.org/arcadia/between-salt-and-water-environmental-crisis-sambhar-lake-rajasthan-india}

  \item The Hans India (2022, April). Sand mafia digging deep into salt lands at
        Chinaganjam.
        \url{https://www.thehansindia.com/andhra-pradesh/sand-mafia-digging-deep-into-salt-lands-at-chinaganjam-736019}

  \item Tata Group / Tata Chemicals Ltd (n.d.). Chemicals With A Capital Sea.
        \url{https://www.tata.com/newsroom/tata-chemicals-with-a-capital-sea}

  \item Sahapedia (n.d.). In Search of White Gold: Salt Harvesting at Marakkanam.
        \url{https://www.sahapedia.org/search-white-gold-salt-harvesting-marakkanam}

  \item GeoJuris Today (2024). Salt Workers Welfare Bill --- Legislative Status and
        Analysis. \url{https://geojuristoday.in/}

  \item NRSC/ISRO (n.d.). \textit{Salt Pan Atlas of India} (Cartosat-2/3).
        \url{https://www.nrsc.gov.in/nrscnew/resources_atlas_SaltPan.php}

  \item Conservation Action Trust (n.d.). How much of Mumbai's salt pans can be
        developed? \url{https://cat.org.in/}

  \item Bombay Natural History Society (BNHS) (2022). Flamingo Count Report, Thane
        Creek. BNHS, Mumbai.

  \item Gupta, A.\ \& Srinivas, S.\ (2026). \textit{Lavana GIS Dataset v1: Indian Salt
        Pan Geodatabase}. Lavana Project Research Initiative.
        \url{https://github.com/ArnavGupta-codes/Lavana-Project-}

  \item Self Employed Women's Association (SEWA) (2024). Solar Pump Initiative for
        Gujarat Salt Workers. \url{https://www.sewa.org/}

  \item Food Safety and Standards Authority of India (FSSAI). IS 7224:2018 ---
        Iodised Salt Specification. Government of India.

  \item CSIR-CSMCRI (Central Salt and Marine Chemicals Research Institute). R\&D on
        Geo-Membrane Liners and Solar-Assisted Evaporation. Bhavnagar, Gujarat.
\end{enumerate}

\end{document}
